\section*{Chapitre \ref{ch:s1:pairs-and-baskets} --- Paires et paniers}

\begin{solution}{exo:s1:pairs-and-baskets:1}
$10 \times \binom{45}{2} = 10 \times 990 = 9\,900$ paires~; un test valide à 5~\% en rejette environ $9\,900 \times 0.05 = 495$, toutes fausses.
\end{solution}

\begin{solution}{exo:s1:pairs-and-baskets:2}
La procédure rejette les $k$ plus petites, où la $k$-ième est au plus $qk/m$~: $0.05 \times 20/1000 = 0.001$.
\end{solution}

\begin{solution}{exo:s1:pairs-and-baskets:3}
L'ouverture négocie 1~\$ sur chaque jambe, et la clôture les négocie à nouveau~: quatre unités à cinq points de base, soit 20 points de base sur une jambe.
\end{solution}

\begin{solution}{exo:s1:pairs-and-baskets:4}
Soit $2 \times 2.4~\% = 4.8~\%$. La règle parie sur la convergence et ne prévoit aucun stop~: la prime d'OPA a fait monter la jambe courte d'un quart en une journée, et le titre
ne s'est plus négocié, de sorte que le spread ne pouvait jamais revenir. La perte vient de la prime, et non du seuil.
\end{solution}

\begin{solution}{exo:s1:pairs-and-baskets:5}
La somme des carrés des écarts est la plus petite pour les titres dont les cours bougent le moins, quelle que soit leur relation. Là où existent des substituts, ceux-ci remportent le classement (311
sur 320 sélections)~; là où il n'en existe aucun, le classement se remplit de titres calmes dont les trajectoires sont restées proches par hasard, et leurs spreads s'écartent comme toute
différence de marches aléatoires.
\end{solution}

\begin{solution}{exo:s1:pairs-and-baskets:6}
La procédure contrôle le taux de fausses découvertes lorsque les p-valeurs des hypothèses nulles vraies sont uniformes (ou conservatrices). Ici ce n'est pas le cas~: les queues épaisses et
la volatilité en grappes rendent les plus petites p-valeurs des 745\,000 paires nulles bien plus petites que ne le donnerait une loi uniforme, et ces fausses paires remplissent la
liste des découvertes devant les paires implantées. La procédure est juste~; son entrée est fausse.
\end{solution}

\begin{solution}{exo:s1:pairs-and-baskets:7}
Parmi les 4\,020 paires rejetées de nouveau, 401 sont implantées, soit 10.0~\%~; parmi les 58\,024 rejets de l'année de formation, la part des paires implantées est de 0.91~\%. Exiger la persistance
multiplie par onze la part des paires implantées~: une paire qui reste cointégrée est bien plus susceptible d'être réelle, ce qui fait l'utilité du résultat nommé.
\end{solution}

\begin{solution}{exo:s1:pairs-and-baskets:8}
La correction pour fausses découvertes suppose des p-valeurs calibrées, ce que le test d'Engle--Granger ne fournit pas sur des rendements à queues épaisses et hétéroscédastiques~; les
survivantes sont surtout fausses. Et sélectionner à l'aide des mêmes données que celles qui ont servi à publier le ratio de Sharpe le surestime. Vérifier la calibration du test sur des données sans
cointégration, exiger la persistance ou un lien économique, et rendre compte sur une période ultérieure.
\end{solution}

\begin{solution}{pb:s1:pairs-and-baskets:1}
\begin{enumerate}
\item Acheter le retardataire et vendre à découvert le meneur de deux titres ayant évolué ensemble, position clôturée quand les deux cours convergent~; la \omterm{def:s1:pairs-and-baskets:distance}{méthode des distances} sélectionne les paires selon la somme des
carrés des écarts de prix normalisés sur une \omterm{def:s1:pairs-and-baskets:distance}{période de formation}, puis les négocie sur la \omterm{def:s1:pairs-and-baskets:distance}{période de trading} suivante.
\item Formation de douze mois, trading de six mois, ouverture à deux écarts-types de formation, clôture au croisement suivant, 20 premières paires~; jusqu'à 11~\% par an de
rendements excédentaires, 1962--2002.
\item Cointégration~: un test d'Engle--Granger sur les logarithmes des cours, en négociant le résidu~; copule~: une copule ajustée aux rendements quotidiens, en négociant un indice de mouvements
conjoints improbables.
\item 91, 85 et 43 points de base par mois avant frais pour la \omterm{def:s1:pairs-and-baskets:distance}{méthode des distances}, la cointégration et la copule~; 38, 33 et 5 après.
\item Cinquante paires de substituts proches dont le spread des logarithmes des cours est un AR(1) stationnaire (dix jours, 2~\%)~; sans elles, aucune paire du marché synthétique n'est réelle.
\item Distances et copule~: 311 sur 320~; deux étapes~: 113~; Engle--Granger~: 40.
\item Marché jumeau~: 5.05, 3.89, 1.03, $-0.17$, $-0.49$~; marché sans substituts~: $-0.43$, $-0.50$, $-0.11$, $-0.47$, $-0.35$.
\item Substituts implantés sélectionnés par distance, ouverts deux fois~; au jour 111, la jambe courte a fait l'objet d'une OPA, le spread a bondi à 35~\% et la paire a perdu 24.6~\% par
dollar de jambe, sans clôture.
\item 745\,465 tests sur seize cycles~; à un seuil nominal de 5~\%, ils rejettent 7.4~\% des paires dont aucune n'est réelle, et à 0.1~\%, trois fois le taux nominal.
\item 97.2~\%~: des p-valeurs mal calibrées placent les fausses paires devant les paires implantées.
\item Il classe les 200 paires les plus proches par p-valeur, et au sein de celles-ci le test privilégie encore ses faux rejets~: 113 paires implantées, pour un net de 1.03.
\item Les mouvements propres d'un titre n'ont aucun pendant chez ses pairs, si bien que le spread contre tout panier est une marche aléatoire.
\item \textbf{Résultat nommé.} 7.9~\% des paires rejetées à 5~\% pendant une année de formation le sont de nouveau l'année suivante (84.6~\% des paires implantées, 7.2~\% des
autres)~; la stratégie des distances dégage un ratio de Sharpe de 5.05 sur le marché avec substituts et de $-0.43$ sur le marché sans.
\item Les paires réelles persistent, les fausses au taux approximatif de faux rejets du test, si bien que la persistance les sépare.
\item Une baisse continue de la rentabilité, avec de bonnes performances pendant les turbulences prolongées, notamment la crise de 2008.
\item Résidus face aux modèles factoriels, valeur relative structurelle avec liens contractuels, et larges portefeuilles de petits paris.
\item Lancer le test sur des données où l'hypothèse nulle est vraie (paires mélangées, ou une simulation reproduisant les queues et la volatilité des données) et comparer les taux de rejet aux
niveaux nominaux.
\item Une raison économique d'être substitut, un spread présentant un retour à la moyenne sur des données non utilisées pour le sélectionner, et de la liquidité sur les deux jambes.
\item Filtrer les OPA et les rumeurs, fixer un cap à la perte par paire, clôturer les paires dont une jambe est sous offre, et diversifier sur de nombreuses paires.
\item Substituts proches dont les cours sont ramenés l'un vers l'autre par autre chose que leur propre passé.
\end{enumerate}
\end{solution}

\begin{solution}{iq:s1:pairs-and-baskets:1}
Partir d'un lien économique (même activité, classes d'actions, un fournisseur et un client), puis vérifier le lien dans les données~: \omterm{def:s1:pairs-and-baskets:distance}{méthode des distances} ou test de cointégration sur
une \omterm{def:s1:pairs-and-baskets:distance}{période de formation}, persistance sur une période ultérieure, liquidité et emprunt de titres sur les deux jambes. Compter chaque paire testée et réserver une période hors échantillon à l'évaluation.
\end{solution}

\begin{solution}{iq:s1:pairs-and-baskets:2}
La corrélation porte sur des rendements qui bougent ensemble jour après jour~; la cointégration porte sur des cours qui restent ensemble dans le temps, soit un spread stationnaire. Deux titres peuvent être
très corrélés et s'écarter l'un de l'autre, ou faiblement corrélés d'un jour à l'autre et pourtant cointégrés. Le \omterm{def:s1:pairs-and-baskets:pairs}{pairs trading} parie sur le spread, si bien que la cointégration compte,
et elle est plus difficile à établir.
\end{solution}

\begin{solution}{iq:s1:pairs-and-baskets:3}
Benjamini--Hochberg ou une borne sur l'erreur familiale (family-wise) sur l'ensemble des tests. Ce qui peut encore mal tourner~: des p-valeurs mal calibrées à cause des queues épaisses et de la volatilité changeante (les vérifier sur
des données nulles)~; la dépendance entre les tests~; et les relations réelles dans l'échantillon qui se rompent ensuite. Exiger une persistance hors échantillon.
\end{solution}

\begin{solution}{iq:s1:pairs-and-baskets:4}
L'hypothèse de la paire n'existe plus~: la clôturer, en prenant la perte si la cible est la jambe courte, sauf si le portefeuille fait délibérément de l'arbitrage de fusions (chapitre~11). Puis
vérifier si les autres paires sont exposées à la même opération ou au même secteur.
\end{solution}

\begin{solution}{iq:s1:pairs-and-baskets:5}
Avec $N$ la matrice $T \times n$ des cours normalisés, $D_{ij} = \|N_i\|^2 + \|N_j\|^2 - 2 N_i^\top N_j$~: un seul produit matriciel $N^\top N$ ($3\,000 \times
3\,000$ à partir de $T \times 3\,000$), puis le triangle supérieur.
\end{solution}

\begin{solution}{iq:s1:pairs-and-baskets:6}
L'écart attendu décroît comme $\phi^h$, si bien qu'il est divisé par deux après $h = \ln 2/(-\ln\phi)$ jours. La statistique de Dickey--Fuller est invariante par changement d'échelle~: sa puissance
dépend de $\phi$ et de la longueur de l'échantillon, non de $\sigma$. Un spread étroit n'est pas plus facile à détecter, seulement moins coûteux à négocier au regard de ses mouvements.
\end{solution}
